% --- CORRECCIÓN IMPORTANTE EN LA PRIMERA LÍNEA ---
% Agregamos "xcolor=table" para que funcione \rowcolor sin errores
\documentclass[aspectratio=169, 10pt, xcolor=table]{beamer}

% --- Configuración del Tema (estilo profesional como el ejemplo) ---
\usetheme{Madrid}
\usecolortheme{dolphin}

% --- Paquetes necesarios ---
\usepackage[utf8]{inputenc}
\usepackage[spanish]{babel}
\usepackage{amsmath, amssymb}
\usepackage{booktabs}
\usepackage{graphicx}
\usepackage{listings}
\usepackage{tikz}
\usetikzlibrary{arrows.meta, positioning, shapes.geometric}

% --- Ruta de Imágenes (crea la carpeta "imagenes" y coloca ahí tus archivos) ---
\graphicspath{{./imagenes/}}

% --- Estilo para código Python ---
\definecolor{codegreen}{rgb}{0,0.6,0}
\definecolor{codegray}{rgb}{0.5,0.5,0.5}
\definecolor{codepurple}{rgb}{0.58,0,0.82}
\definecolor{backcolour}{rgb}{0.95,0.95,0.92}
\lstset{
    language=Python,
    backgroundcolor=\color{backcolour},
    commentstyle=\color{codegreen},
    keywordstyle=\color{blue},
    numberstyle=\tiny\color{codegray},
    stringstyle=\color{codepurple},
    basicstyle=\ttfamily\scriptsize,
    breaklines=true,
    frame=single,
    showstringspaces=false
}

% --- Pie de página personalizado ---
\makeatletter

\setbeamertemplate{footline}{
  \leavevmode%
  \hbox{%
  \begin{beamercolorbox}[wd=.333333\paperwidth,ht=2.25ex,dp=1ex,center]{author in head/foot}%
    \usebeamerfont{author in head/foot}\insertshortauthor
  \end{beamercolorbox}%
  \begin{beamercolorbox}[wd=.333333\paperwidth,ht=2.25ex,dp=1ex,center]{title in head/foot}%
    \usebeamerfont{title in head/foot}Sesión 10
  \end{beamercolorbox}%
  \begin{beamercolorbox}[wd=.333333\paperwidth,ht=2.25ex,dp=1ex,right]{date in head/foot}%
    \usebeamerfont{date in head/foot}Machine Learning \hfill \insertframenumber/\inserttotalframenumber \hspace*{0.3cm}
  \end{beamercolorbox}}%
  \vskip0pt%
}

\makeatother

% --- Datos de la portada ---
\title[SESIÓN 10]{\textbf{Sesión 10}}
\subtitle{Series de Tiempo: \\ Modelado estadístico y aprendizaje automático}
\author[Prof. C. Sánchez]{Carlos César Sánchez Coronel}
\institute{\textbf{MACHINE LEARNING}}
\date{}

% Logo opcional (descomentar si tienes la imagen)
% \titlegraphic{\includegraphics[height=1.5cm]{logo.png}}

% --- Eliminar la barra de navegación (opcional) ---
\setbeamertemplate{navigation symbols}{}

% --- Índice al inicio de cada sección ---
\AtBeginSection[]{
  \begin{frame}
    \frametitle{Contenido}
    \tableofcontents[currentsection]
  \end{frame}
}

% ============================================================================
% INICIO DEL DOCUMENTO
% ============================================================================
\begin{document}

% --- Portada ---
\begin{frame}
    \titlepage
\end{frame}

% --- Índice general ---
\begin{frame}{Contenido}
    \tableofcontents
\end{frame}

% ============================================================================
% SECCIÓN 1: INTRODUCCIÓN Y COMPONENTES
% ============================================================================
\section{Componentes de una Serie Temporal}

\begin{frame}{¿Qué es una serie temporal?}
    \begin{itemize}
        \item Datos observados en secuencia temporal (diario, mensual, etc.).
        \item Ejemplos: ventas diarias, tráfico web, indicadores económicos.
        \item Objetivo: entender patrones y pronosticar valores futuros.
    \end{itemize}
    \vspace{0.3cm}
    \begin{block}{Componentes fundamentales}
        \begin{description}
            \item[Tendencia:] movimiento a largo plazo (creciente, decreciente, estable).
            \item[Estacionalidad:] patrones regulares en períodos fijos (diario, semanal, anual).
            \item[Ciclo:] fluctuaciones sin período fijo (asociadas a condiciones económicas).
            \item[Ruido:] variaciones aleatorias no explicadas.
        \end{description}
    \end{block}
    % Basado en Pregunta 1
    % IMAGEN: Serie con tendencia y estacionalidad
\end{frame}

\begin{frame}{Descomposición de series}
    \begin{columns}[t]
        \column{0.5\textwidth}
        \textbf{Modelo aditivo}
        \[
        y_t = T_t + S_t + C_t + R_t
        \]
        Amplitud estacional constante.
        \column{0.5\textwidth}
        \textbf{Modelo multiplicativo}
        \[
        y_t = T_t \times S_t \times C_t \times R_t
        \]
        Amplitud estacional proporcional al nivel.
    \end{columns}
    \vspace{0.3cm}
    \begin{itemize}
        \item Método de descomposición: \textbf{STL} (Seasonal and Trend decomposition using Loess).
        \item Ajuste estacional: eliminar la estacionalidad para analizar tendencia y ciclo.
    \end{itemize}
    % Basado en Pregunta 2
\end{frame}

% ============================================================================
% SECCIÓN 2: FEATURE ENGINEERING PARA SERIES
% ============================================================================
\section{Feature Engineering para Series Temporales}

\begin{frame}{Creación de características supervisadas}
    \begin{itemize}
        \item Convertir el problema de series en aprendizaje supervisado generando predictores a partir de la historia.
    \end{itemize}
    \begin{columns}[t]
        \column{0.5\textwidth}
        \textbf{Rezagos (lags)}
        \begin{itemize}
            \item $y_{t-1}, y_{t-2}, \dots, y_{t-L}$
            \item Se seleccionan usando ACF/PACF.
        \end{itemize}
        \textbf{Ventanas móviles}
        \begin{itemize}
            \item Media móvil, desviación estándar, mínimo, máximo.
            \item Capturan tendencia local y volatilidad.
        \end{itemize}
        \column{0.5\textwidth}
        \textbf{Diferencias}
        \begin{itemize}
            \item $\Delta y_t = y_t - y_{t-1}$
            \item Para hacer la serie estacionaria.
        \end{itemize}
        \textbf{Variables de calendario}
        \begin{itemize}
            \item Hora, día de semana, mes, año.
            \item Festivos, fines de semana.
        \end{itemize}
    \end{columns}
    % Basado en Pregunta 3
\end{frame}

\begin{frame}{Codificación de variables cíclicas}
    \begin{problem}
        Codificar hora como entero (0-23) no respeta la naturaleza circular (23 y 0 son cercanas).
    \end{problem}
    \begin{block}{Solución: transformación seno/coseno}
        \[
        \sin\left(\frac{2\pi \cdot \text{hora}}{24}\right),\quad \cos\left(\frac{2\pi \cdot \text{hora}}{24}\right)
        \]
        Se aplica también a día de semana (período 7), mes (período 12), etc.
    \end{block}
    % Basado en Pregunta 4
\end{frame}

% ============================================================================
% SECCIÓN 3: MODELOS CLÁSICOS: ARIMA Y SARIMA
% ============================================================================
\section{Modelos Clásicos: ARIMA y SARIMA}

\begin{frame}{ARMA(p,q)}
    \begin{block}{Ecuación}
        \[
        y_t = c + \phi_1 y_{t-1} + \cdots + \phi_p y_{t-p} + \varepsilon_t + \theta_1 \varepsilon_{t-1} + \cdots + \theta_q \varepsilon_{t-q}
        \]
        donde $\varepsilon_t$ es ruido blanco.
    \end{block}
    \begin{itemize}
        \item AR(p): parte autorregresiva.
        \item MA(q): parte de medias móviles.
    \end{itemize}
\end{frame}

\begin{frame}{ARIMA(p,d,q)}
    \begin{itemize}
        \item Para series no estacionarias se aplican diferencias.
        \item $d$: número de diferencias necesarias para lograr estacionariedad.
        \item ARIMA(p,d,q) = ARMA(p,q) aplicado a la serie diferenciada $d$ veces.
    \end{itemize}
    \vspace{0.3cm}
    \begin{block}{Estacionariedad}
        \begin{itemize}
            \item Media y varianza constantes en el tiempo.
            \item Se verifica con test de Dickey-Fuller (ADF).
        \end{itemize}
    \end{block}
    % Basado en Pregunta 5
\end{frame}

\begin{frame}{Identificación de órdenes con ACF y PACF}
    \begin{columns}[t]
        \column{0.5\textwidth}
        \textbf{ACF (autocorrelación)}
        \begin{itemize}
            \item Correlación con rezagos.
            \item Para MA(q): corta después del rezago q.
        \end{itemize}
        \column{0.5\textwidth}
        \textbf{PACF (autocorrelación parcial)}
        \begin{itemize}
            \item Correlación eliminando dependencias intermedias.
            \item Para AR(p): corta después del rezago p.
        \end{itemize}
    \end{columns}
    \vspace{0.3cm}
    \begin{exampleblock}{Ejemplo}
        PACF significativa en rezago 1 y luego cero $\Rightarrow$ AR(1). \\
        ACF significativa en rezago 1 y luego cero $\Rightarrow$ MA(1).
    \end{exampleblock}
    % Basado en Pregunta 6
\end{frame}

\begin{frame}{SARIMA: incorporando estacionalidad}
    \begin{block}{Notación SARIMA(p,d,q)(P,D,Q)$_s$}
        \begin{itemize}
            \item $s$: período estacional (ej. 12 para mensual, 7 para diario).
            \item $P, D, Q$: órdenes AR, diferencias, MA estacionales.
        \end{itemize}
    \end{block}
    \begin{itemize}
        \item Diferencias regulares ($d$) eliminan tendencia.
        \item Diferencias estacionales ($D$) eliminan estacionalidad (ej. $y_t - y_{t-s}$).
        \item Selección con criterios AIC, BIC.
    \end{itemize}
    % Basado en Pregunta 7
\end{frame}

% ============================================================================
% SECCIÓN 4: MODELOS DE MACHINE LEARNING
% ============================================================================
\section{Modelos de Machine Learning}

\begin{frame}{XGBoost para series temporales}
    \begin{itemize}
        \item Requiere crear características (lags, ventanas, variables temporales).
        \item Ventajas: captura no linealidades, múltiples predictores externos.
        \item Desventajas: no extrapola tendencias fuera del rango de entrenamiento.
    \end{itemize}
    \vspace{0.3cm}
    \begin{block}{Características importantes}
        \begin{itemize}
            \item Lags de la variable objetivo.
            \item Medias móviles y desviaciones.
            \item Variables de calendario (día de semana, mes, festivos).
            \item Variables exógenas (clima, promociones).
        \end{itemize}
    \end{block}
    % Basado en Pregunta 8
\end{frame}

\begin{frame}{Prophet (Facebook)}
    \begin{block}{Modelo aditivo}
        \[
        y(t) = g(t) + s(t) + h(t) + \varepsilon_t
        \]
        \begin{itemize}
            \item $g(t)$: tendencia (lineal por tramos o logística).
            \item $s(t)$: estacionalidad (series de Fourier).
            \item $h(t)$: efectos de días festivos.
        \end{itemize}
    \end{block}
    \begin{itemize}
        \item Robusto a datos faltantes y outliers.
        \item Fácil de usar, interpretable.
        \item Permite incluir regresores externos (versión reciente).
    \end{itemize}
    % Basado en Pregunta 9
\end{frame}

\begin{frame}{Comparación de enfoques}
    \begin{table}
        \centering
        \begin{tabular}{l|c|c|c}
            \toprule
            \textbf{Característica} & \textbf{ARIMA/SARIMA} & \textbf{XGBoost} & \textbf{Prophet} \\
            \midrule
            Variables exógenas & Limitado (ARIMAX) & Sí & Sí (regresores) \\
            Estacionalidad & Explícita (SARIMA) & Variables dummy & Fourier \\
            Extrapolación de tendencia & Sí & Limitada & Sí \\
            Interpretabilidad & Media & Baja & Alta \\
            Facilidad de uso & Media & Alta & Muy alta \\
            \bottomrule
        \end{tabular}
    \end{table}
    % Basado en Pregunta 10
\end{frame}

% ============================================================================
% SECCIÓN 5: MÉTRICAS DE EVALUACIÓN
% ============================================================================
\section{Métricas de Evaluación}

\begin{frame}{Métricas para pronósticos}
    \begin{columns}[t]
        \column{0.5\textwidth}
        \textbf{MAE} (Mean Absolute Error)
        \[
        \frac{1}{n}\sum |y_t - \hat{y}_t|
        \]
        \textbf{RMSE} (Root Mean Squared Error)
        \[
        \sqrt{\frac{1}{n}\sum (y_t - \hat{y}_t)^2}
        \]
        \column{0.5\textwidth}
        \textbf{MAPE} (Mean Absolute Percentage Error)
        \[
        \frac{100}{n}\sum \left|\frac{y_t - \hat{y}_t}{y_t}\right|
        \]
        \textbf{MASE} (Mean Absolute Scaled Error)
        \[
        \frac{\frac{1}{n}\sum |y_t - \hat{y}_t|}{\frac{1}{n-1}\sum |y_t - y_{t-1}|}
        \]
    \end{columns}
    \vspace{0.3cm}
    \begin{itemize}
        \item MASE < 1: modelo supera al naive (pronóstico del día anterior).
        \item MAPE tiene problemas con valores cercanos a cero.
    \end{itemize}
    % Basado en Preguntas 11 y 12
\end{frame}

% ============================================================================
% SECCIÓN 6: VALIDACIÓN EN SERIES TEMPORALES
% ============================================================================
\section{Validación en Series Temporales}

\begin{frame}{Time Series Split}
    \begin{alertblock}{¡No usar k-fold estándar!}
        Mezcla datos futuros con pasados, causando data leakage.
    \end{alertblock}
    \begin{columns}[t]
        \column{0.5\textwidth}
        \textbf{Expanding window}
        \begin{itemize}
            \item Entrenamiento creciente.
            \item Validación fija.
        \end{itemize}
        \column{0.5\textwidth}
        \textbf{Rolling window}
        \begin{itemize}
            \item Ventana de tamaño fijo se desliza.
        \end{itemize}
    \end{columns}
    \vspace{0.3cm}
    En scikit-learn: \texttt{TimeSeriesSplit}.
    % Basado en Pregunta 13
\end{frame}

% ============================================================================
% SECCIÓN 7: CASO INTEGRADO: PRONÓSTICO DE VENTAS EN RETAIL
% ============================================================================
\section{Caso Integrado: Pronóstico de Ventas en Retail}

\begin{frame}{Problema y datos}
    \begin{itemize}
        \item Cadena de supermercados: pronosticar ventas diarias por tienda/producto.
        \item Datos: 5 años de ventas, precios, promociones, festivos, clima.
        \item Horizonte: 7 días, actualización diaria.
    \end{itemize}
    % Basado en Pregunta 20
\end{frame}

\begin{frame}{Feature engineering específico}
    \begin{enumerate}
        \item Rezagos: ventas de los últimos 7, 14, 28 días.
        \item Medias móviles: media de últimos 7 días, media del mismo día de la semana en últimas 4 semanas.
        \item Variables de calendario: día de semana, mes, año, fin de semana, festivos.
        \item Clima: temperatura, lluvia (y rezagos).
        \item Promociones: precio actual, indicador de promoción, días desde inicio.
    \end{enumerate}
\end{frame}

\begin{frame}{Comparación de modelos}
    \begin{table}
        \centering
        \begin{tabular}{l|c|c|c}
            \toprule
            \textbf{Modelo} & \textbf{MAE} & \textbf{RMSE} & \textbf{MASE} \\
            \midrule
            Naive (día anterior) & 2450 & 3200 & 1.00 \\
            SARIMA & 1850 & 2500 & 0.75 \\
            XGBoost & 1420 & 1950 & 0.58 \\
            Prophet & 1600 & 2100 & 0.65 \\
            \bottomrule
        \end{tabular}
    \end{table}
    \begin{itemize}
        \item XGBoost obtiene mejor desempeño, pero Prophet es más interpretable.
        \item En producción se implementa XGBoost con reentrenamiento diario.
    \end{itemize}
\end{frame}

% ============================================================================
% SECCIÓN 8: IMPLEMENTACIÓN EN PYTHON
% ============================================================================
\section{Implementación en Python}

\begin{frame}[fragile]{ARIMA con statsmodels}
    \begin{lstlisting}
import statsmodels.api as sm
from statsmodels.tsa.arima.model import ARIMA

model = ARIMA(series, order=(1,1,1), seasonal_order=(1,1,1,7))
fitted = model.fit()
forecast = fitted.forecast(steps=28)
    \end{lstlisting}
\end{frame}

\begin{frame}[fragile]{XGBoost con validación temporal}
    \begin{lstlisting}
import xgboost as xgb
from sklearn.model_selection import TimeSeriesSplit

tscv = TimeSeriesSplit(n_splits=5)
for train_idx, val_idx in tscv.split(X):
    X_train, X_val = X[train_idx], X[val_idx]
    y_train, y_val = y[train_idx], y[val_idx]
    model = xgb.XGBRegressor(n_estimators=100, learning_rate=0.1)
    model.fit(X_train, y_train)
    # evaluar...
    \end{lstlisting}
\end{frame}

\begin{frame}[fragile]{Prophet}
    \begin{lstlisting}
from prophet import Prophet
import pandas as pd

df = pd.DataFrame({'ds': fechas, 'y': ventas})
model = Prophet(yearly_seasonality=True, weekly_seasonality=True)
model.add_country_holidays(country_name='PE')
model.fit(df)
future = model.make_future_dataframe(periods=28)
forecast = model.predict(future)
    \end{lstlisting}
\end{frame}

% ============================================================================
% SECCIÓN 9: CONEXIÓN Y RESUMEN
% ============================================================================
\section{Conexión y Resumen}

\begin{frame}{Conexión con el resto del curso}
    \begin{itemize}
        \item Feature engineering para series usa transformaciones y agregaciones.
        \item XGBoost (Sesión 7) se aplica aquí con características temporales.
        \item Validación cruzada temporal (Sesión 8) es clave.
        \item Métricas de regresión (MAE, RMSE) ya conocidas.
    \end{itemize}
\end{frame}

\begin{frame}{Lo que hemos aprendido}
    \begin{exampleblock}{Conceptos clave}
        \begin{itemize}
            \item Componentes de series: tendencia, estacionalidad, ciclo, ruido.
            \item Descomposición aditiva/multiplicativa, STL.
            \item Feature engineering: lags, ventanas, variables cíclicas, calendario.
            \item Modelos clásicos: ARIMA, SARIMA, identificación con ACF/PACF.
            \item Machine learning: XGBoost, Prophet.
            \item Métricas: MAE, RMSE, MAPE, MASE.
            \item Validación temporal: TimeSeriesSplit.
            \item Caso práctico: pronóstico de ventas en retail.
        \end{itemize}
    \end{exampleblock}
\end{frame}

% --- Diapositiva final ---
\begin{frame}
    \centering
    \Huge \textbf{¡Gracias!}

\end{frame}

\end{document}